\subsection{Loads Along the Nominal Descent}
\label{sec:flight-exposure}
A safe descent depends on the loads during the whole flight, not only at touchdown. Table~\ref{tab:flight-exposure} uses the same 25.97 t spherical trajectory as the propellant answer, so loads and budget refer to one design. Dynamic pressure is $q=\rho U^2/2$, where $U$ is the speed relative to the air~\cite{dynamicPressure}. The specific force $\lVert\mathbf F_a+\mathbf T\rVert/(m g_0)$ is the non-gravitational acceleration that the structure actually carries.

\begin{table}[htbp]
\centering
\caption{Peak loads on the spherical reference trajectory, from the replay at a 0.0025 s step. Height is the radial altitude of the center of mass. The planar design gave 297.3 kPa, $7.70\,g_0$, and 9.62 MN.}
\label{tab:flight-exposure}
\begin{tabular}{@{}lrrl@{}}
\toprule
Quantity & Time (s) & Height (km) & Peak \\
\midrule
Dynamic pressure & 175.725 & 9.834 & 292.9 kPa \\
Specific-force magnitude & 187.413 & 3.343 & $7.84\,g_0$ \\
Normal aerodynamic force & 187.523 & 3.302 & 9.98 MN \\
\bottomrule
\end{tabular}
\end{table}

Dynamic pressure peaks before the engine lights at 181.32 s, while the specific force peaks during the burn. Near that second peak the body axis is about $20^\circ$ from the direction of the oncoming air, which produces a large side force. The attitude therefore matters as much as the speed and the air density. Because structural and thermal capacity are not modeled, the table states what a booster flying this trajectory would have to withstand; it does not show that it can.

\subsection{Aerodynamic Heating Estimate}
\label{sec:heating}
To gauge the heating, we apply the Sutton--Graves correlation for the convective heat flux at a stagnation point in air~\cite{suttonGraves},
\begin{equation}\label{eq:sutton}
 \dot q_s=k\sqrt{\frac{\rho}{R_n}}\,U^3,\qquad k=1.7415\times10^{-4}\ \mathrm{kg^{1/2}\,m^{-1}},
\end{equation}
along the reference trajectory, where $R_n$ is the effective nose radius of the windward surface. The booster changes attitude during the descent and meets the air partly tail first, so $R_n$ is uncertain. We bracket it between 1 m, for a rounded edge, and the 3.5 m body radius. The peak occurs at 168.7 s, at 15.5 km altitude and 1500 m/s, before the dynamic-pressure peak. It is 249 kW/m$^2$ for $R_n=1$ m and 133 kW/m$^2$ for $R_n=3.5$ m, and the heat absorbed per unit area over the whole flight is 8.6 and 4.6 MJ/m$^2$. Equation~\eqref{eq:sutton} is intended for blunt bodies at hypersonic speed, and the separation speed is only about Mach 5, so these values indicate the order of magnitude. They are the design heating that a thermal protection system would have to handle; sizing that system is outside this study.

\subsection{Cost of Limiting Dynamic Pressure}
\label{sec:qcap}
The 293 kPa peak occurs while the engine is off, so a later correction cannot reduce it. To see what a lower structural limit would cost, we add the constraint $q\le q_{max}$ to the spherical optimization at every mesh point, with a 1\% margin. Starting trajectories come from the feedback search with one extra braking burn high in the descent, followed by a second coast and the landing burn. A solution is accepted only if both replays meet every touchdown limit and the replayed peak stays at or below $q_{max}$. We have no sourced load limit for this vehicle, so the 200 and 150 kPa caps are examples.

\begin{table}[htbp]
\centering
\caption{Dynamic-pressure caps for the $15^\circ$ spherical case. Each row is the best accepted solution for that limit. Masses in tonnes.}
\label{tab:qcap}
\begin{tabular}{@{}lllrrrr@{}}
\toprule
Cap (kPa) & Status & Burns & Inventory & Entry burn & Landing burn & Peak $q$ (kPa) \\
\midrule
none & optimized & 1 & 25.97 & -- & 20.89 & 292.9 \\
200 & optimized & 2 & 37.30 & 13.50 & 18.72 & 198.0 \\
150 & optimized & 2 & 50.17 & 27.08 & 18.05 & 148.7 \\
\bottomrule
\end{tabular}
\end{table}

A 200 kPa cap raises the separation inventory from 25.97 t to 37.30 t. The added propellant goes into a 38 s entry burn starting at 43.2 km altitude, while the landing burn still uses about 19 t. A 150 kPa cap needs 50.17 t, with a 77 s entry burn starting at 62.7 km altitude. Lowering the limit from 200 to 150 kPa costs another 12.9 t. These solutions pass the same replay checks as the reference. They were still improving when the optimization was stopped, so they are upper estimates of the minimum.


At the masses involved, the 2.846 MN engine decelerates the booster at roughly $1.5$--$2\,g_0$, far less than aerodynamic drag near peak dynamic pressure. Lowering $q$ therefore takes a long burn while the booster is still fast, and all the propellant burned early has to be carried from separation. If a structural limit near or below 200 kPa applied, it rather than the landing burn would set the propellant requirement.
